% ============================================================
% UDE THESIS PREAMBLE
% University of Duisburg-Essen | Chair of Econometrics
% ============================================================
% Contains: document class, packages, layout, metadata, and
%           acronym definitions.
%
% Included by ude_thesis.tex via: % ============================================================
% UDE THESIS PREAMBLE
% University of Duisburg-Essen | Chair of Econometrics
% ============================================================
% Contains: document class, packages, layout, metadata, and
%           acronym definitions.
%
% Included by ude_thesis.tex via: % ============================================================
% UDE THESIS PREAMBLE
% University of Duisburg-Essen | Chair of Econometrics
% ============================================================
% Contains: document class, packages, layout, metadata, and
%           acronym definitions.
%
% Included by ude_thesis.tex via: % ============================================================
% UDE THESIS PREAMBLE
% University of Duisburg-Essen | Chair of Econometrics
% ============================================================
% Contains: document class, packages, layout, metadata, and
%           acronym definitions.
%
% Included by ude_thesis.tex via: \input{ude_thesis_preamble}
% Do NOT compile this file directly — compile ude_thesis.tex.
% ============================================================

\documentclass[11pt,a4paper]{article}


% ============================================================
% LANGUAGE — change this one value to switch the whole document
% Options: english | german
% ============================================================
\newcommand{\thesislanguage}{english}  % <-- Set language here: english | german

% Include statutory declaration page at end of document?
% true  = declaration page printed (required for final submission)
% false = declaration page omitted  (useful for drafts / seminar papers)
\newcommand{\includestatutorydeclaration}{true}  % <-- true | false


% ============================================================
% FONTS & ENCODING
% ============================================================

\usepackage{lmodern}          % Latin Modern font (scalable)
\usepackage{amssymb,amsmath}  % AMS math symbols and environments
\usepackage{ifxetex,ifluatex} % Detect compiler (pdflatex vs xetex/luatex)
\usepackage{fixltx2e}         % Provides \textsubscript (legacy fix)

% Font encoding: T1 for pdflatex, fontspec for xetex/luatex
\ifnum 0\ifxetex 1\fi\ifluatex 1\fi=0
  % pdflatex path
  \usepackage[T1]{fontenc}
  \usepackage[utf8]{inputenc}
\else
  % xetex / luatex path
  \ifxetex
    \usepackage{mathspec}
    \usepackage{xltxtra,xunicode}
  \else
    \usepackage{fontspec}
  \fi
  \defaultfontfeatures{Mapping=tex-text,Scale=MatchLowercase}
  \newcommand{\euro}{€}
\fi

% Straight quotes in verbatim environments (if upquote.sty is available)
\IfFileExists{upquote.sty}{\usepackage{upquote}}{}


% ============================================================
% PAGE GEOMETRY
% ============================================================

% Wide left margin (5 cm) for binding; standard margins elsewhere
\usepackage[lmargin=5cm,rmargin=2.5cm,tmargin=2.5cm,bmargin=2.5cm]{geometry}


% ============================================================
% FIGURE PLACEMENT
% ============================================================

% Force all figures to appear at their exact position in the source [H]
\usepackage{float}
\let\origfigure\figure
\let\endorigfigure\endfigure
\renewenvironment{figure}[1][2]{
    \expandafter\origfigure\expandafter[H]
}{
    \endorigfigure
}


% ============================================================
% BIBLIOGRAPHY
% ============================================================

% APA citation style via biblatex + biber backend
% maxbibnames=99 ensures all authors are listed in the reference list
\usepackage{csquotes}
\usepackage[backend=biber, maxbibnames=99, style=apa]{biblatex}
\setlength\bibitemsep{1.5\itemsep}         % Extra spacing between reference entries
\bibliography{Resources/references.bib}    % Path to the .bib file


% ============================================================
% CODE SYNTAX HIGHLIGHTING
% ============================================================
% Required for displaying formatted code blocks.
% Safe to remove this entire section if your thesis contains no code.

\usepackage{color}
\usepackage{fancyvrb}
\newcommand{\VerbBar}{|}
\newcommand{\VERB}{\Verb[commandchars=\\\{\}]}
\DefineVerbatimEnvironment{Highlighting}{Verbatim}{commandchars=\\\{\}}
\usepackage{framed}
\definecolor{shadecolor}{RGB}{248,248,248}
\newenvironment{Shaded}{\begin{snugshade}}{\end{snugshade}}
% Token color definitions for syntax highlighting
\newcommand{\AlertTok}[1]{\textcolor[rgb]{0.94,0.16,0.16}{#1}}
\newcommand{\AnnotationTok}[1]{\textcolor[rgb]{0.56,0.35,0.01}{\textbf{\textit{#1}}}}
\newcommand{\AttributeTok}[1]{\textcolor[rgb]{0.13,0.29,0.53}{#1}}
\newcommand{\BaseNTok}[1]{\textcolor[rgb]{0.00,0.00,0.81}{#1}}
\newcommand{\BuiltInTok}[1]{#1}
\newcommand{\CharTok}[1]{\textcolor[rgb]{0.31,0.60,0.02}{#1}}
\newcommand{\CommentTok}[1]{\textcolor[rgb]{0.56,0.35,0.01}{\textit{#1}}}
\newcommand{\CommentVarTok}[1]{\textcolor[rgb]{0.56,0.35,0.01}{\textbf{\textit{#1}}}}
\newcommand{\ConstantTok}[1]{\textcolor[rgb]{0.56,0.35,0.01}{#1}}
\newcommand{\ControlFlowTok}[1]{\textcolor[rgb]{0.13,0.29,0.53}{\textbf{#1}}}
\newcommand{\DataTypeTok}[1]{\textcolor[rgb]{0.13,0.29,0.53}{#1}}
\newcommand{\DecValTok}[1]{\textcolor[rgb]{0.00,0.00,0.81}{#1}}
\newcommand{\DocumentationTok}[1]{\textcolor[rgb]{0.56,0.35,0.01}{\textbf{\textit{#1}}}}
\newcommand{\ErrorTok}[1]{\textcolor[rgb]{0.64,0.00,0.00}{\textbf{#1}}}
\newcommand{\ExtensionTok}[1]{#1}
\newcommand{\FloatTok}[1]{\textcolor[rgb]{0.00,0.00,0.81}{#1}}
\newcommand{\FunctionTok}[1]{\textcolor[rgb]{0.13,0.29,0.53}{\textbf{#1}}}
\newcommand{\ImportTok}[1]{#1}
\newcommand{\InformationTok}[1]{\textcolor[rgb]{0.56,0.35,0.01}{\textbf{\textit{#1}}}}
\newcommand{\KeywordTok}[1]{\textcolor[rgb]{0.13,0.29,0.53}{\textbf{#1}}}
\newcommand{\NormalTok}[1]{#1}
\newcommand{\OperatorTok}[1]{\textcolor[rgb]{0.81,0.36,0.00}{\textbf{#1}}}
\newcommand{\OtherTok}[1]{\textcolor[rgb]{0.56,0.35,0.01}{#1}}
\newcommand{\PreprocessorTok}[1]{\textcolor[rgb]{0.56,0.35,0.01}{\textit{#1}}}
\newcommand{\RegionMarkerTok}[1]{#1}
\newcommand{\SpecialCharTok}[1]{\textcolor[rgb]{0.81,0.36,0.00}{\textbf{#1}}}
\newcommand{\SpecialStringTok}[1]{\textcolor[rgb]{0.31,0.60,0.02}{#1}}
\newcommand{\StringTok}[1]{\textcolor[rgb]{0.31,0.60,0.02}{#1}}
\newcommand{\VariableTok}[1]{\textcolor[rgb]{0.00,0.00,0.00}{#1}}
\newcommand{\VerbatimStringTok}[1]{\textcolor[rgb]{0.31,0.60,0.02}{#1}}
\newcommand{\WarningTok}[1]{\textcolor[rgb]{0.56,0.35,0.01}{\textbf{\textit{#1}}}}


% ============================================================
% TABLES & GRAPHICS
% ============================================================

\usepackage{longtable,booktabs}  % Multi-page tables and publication-quality rules
\usepackage{graphicx}

% Auto-scale images so they never exceed the text area
\makeatletter
\def\maxwidth{\ifdim\Gin@nat@width>\linewidth\linewidth\else\Gin@nat@width\fi}
\def\maxheight{\ifdim\Gin@nat@height>\textheight\textheight\else\Gin@nat@height\fi}
\makeatother
\setkeys{Gin}{width=\maxwidth,height=\maxheight,keepaspectratio}

% Passthrough wrapper used by figure includes; keeps natural image size
\newcommand{\pandocbounded}[1]{#1}


% ============================================================
% HYPERLINKS & PDF METADATA
% ============================================================

% Update pdfauthor and pdftitle to match your thesis
\ifxetex
  \usepackage[setpagesize=false, % page size defined by xetex
              unicode=false,     % unicode breaks when used with xetex
              xetex]{hyperref}
\else
  \usepackage[unicode=true]{hyperref}
\fi
\hypersetup{
  breaklinks  = true,
  bookmarks   = true,
  pdfauthor   = {Name},   % <-- Your name
  pdftitle    = {Title},  % <-- Your thesis title
  colorlinks  = true,
  citecolor   = blue,
  urlcolor    = blue,
  linkcolor   = black,
  pdfborder   = {0 0 0}
}
\urlstyle{same}  % URLs use the same font as surrounding text


% ============================================================
% PARAGRAPH & SECTION FORMATTING
% ============================================================

\setlength{\parindent}{0pt}           % No paragraph indentation
\setlength{\parskip}{6pt plus 2pt minus 1pt}  % Vertical space between paragraphs
\setlength{\emergencystretch}{3em}    % Prevents overfull lines
\setcounter{secnumdepth}{5}           % Number sections down to subsubparagraph level


% ============================================================
% TITLE PAGE COMMANDS
% ============================================================

\usepackage{titling}

% \subtitle{} command — used on the title page
\newcommand{\subtitle}[1]{
  \posttitle{
    \begin{center}\large#1\end{center}
  }
}

% Title metadata — update these for your thesis
\setlength{\droptitle}{-2em}
  \title{Title}                          % <-- Thesis title
  \pretitle{\vspace{\droptitle}\centering\huge}
  \posttitle{\par}
\subtitle{Subtitle}                      % <-- Thesis subtitle
  \author{Name}                          % <-- Your name
  \preauthor{\centering\large\emph}
  \postauthor{\par}
  \predate{\centering\large\emph}
  \postdate{\par}
  \date{today}                           % <-- Submission date


% ============================================================
% ACRONYMS / LIST OF ABBREVIATIONS
% ============================================================

% glossaries-package in acronym-only mode (no page numbers in list)
\usepackage[acronym,nomain,nonumberlist]{glossaries}
\usepackage{supertabular}
\setacronymstyle{short-long}  % First use: SHORT (Long Form); later: SHORT
\makenoidxglossaries

% Styling: bold acronym names, no blank lines between groups
\renewcommand{\glsnamefont}[1]{\textbf{#1}}
\renewcommand{\glsgroupskip}{}
\setlength{\glsdescwidth}{0.8\linewidth}

% Two-column table style: "Abbreviation | Definition" with header and rule
\newglossarystyle{mysuper}{%
  \setglossarystyle{super}
  \renewcommand{\glossentry}[2]{%
    \glsentryitem{##1}\glstarget{##1}{\glossentryname{##1}} &
    \glossentrydesc{##1}\glspostdescription\space ##2\\[0.75ex]
  }%
  \renewcommand{\superheader}{\textbf{Abbreviation} & \textbf{Definition}\\\hline\endhead}%
}

% Acronym definitions — add your own with: \newacronym{label}{SHORT}{Long Form}
\newacronym{itt}{ITT}{Intention-to-Treat}
\newacronym{iv}{IV}{Instrumental Variables}
\newacronym{sbcf}{SBCF}{Shrinkage Bayesian Causal Forest}
\newacronym{sbcf-iv}{SBCF-IV}{Shrinkage Bayesian Instrumental Variable Causal Forest}


% ============================================================
% LINE SPACING
% ============================================================

\usepackage{setspace}
\onehalfspacing   % 1.5 line spacing (standard requirement for academic work)


% ============================================================
% LANGUAGE — babel loaded based on \thesislanguage (set at top)
% ============================================================
% \IfLanguageName{english}{...}{...} in ude_thesis.tex body
% automatically switches text once babel is loaded here.

\usepackage{ifthen}
\usepackage{iflang}
\usepackage[super]{nth}  % Ordinal suffixes: 1st, 2nd, 3rd ...

\ifthenelse{\equal{\thesislanguage}{german}}{
  \usepackage[ngerman]{babel}
}{
  \usepackage[english]{babel}
}


% ============================================================
% MICROTYPOGRAPHY
% ============================================================

% Improves text justification and glyph spacing where available
\IfFileExists{microtype.sty}{%
  \usepackage[nopatch=footnote]{microtype}
  \UseMicrotypeSet[protrusion]{basicmath}  % Disable protrusion for tt fonts
}{}

% Do NOT compile this file directly — compile ude_thesis.tex.
% ============================================================

\documentclass[11pt,a4paper]{article}


% ============================================================
% LANGUAGE — change this one value to switch the whole document
% Options: english | german
% ============================================================
\newcommand{\thesislanguage}{english}  % <-- Set language here: english | german

% Include statutory declaration page at end of document?
% true  = declaration page printed (required for final submission)
% false = declaration page omitted  (useful for drafts / seminar papers)
\newcommand{\includestatutorydeclaration}{true}  % <-- true | false


% ============================================================
% FONTS & ENCODING
% ============================================================

\usepackage{lmodern}          % Latin Modern font (scalable)
\usepackage{amssymb,amsmath}  % AMS math symbols and environments
\usepackage{ifxetex,ifluatex} % Detect compiler (pdflatex vs xetex/luatex)
\usepackage{fixltx2e}         % Provides \textsubscript (legacy fix)

% Font encoding: T1 for pdflatex, fontspec for xetex/luatex
\ifnum 0\ifxetex 1\fi\ifluatex 1\fi=0
  % pdflatex path
  \usepackage[T1]{fontenc}
  \usepackage[utf8]{inputenc}
\else
  % xetex / luatex path
  \ifxetex
    \usepackage{mathspec}
    \usepackage{xltxtra,xunicode}
  \else
    \usepackage{fontspec}
  \fi
  \defaultfontfeatures{Mapping=tex-text,Scale=MatchLowercase}
  \newcommand{\euro}{€}
\fi

% Straight quotes in verbatim environments (if upquote.sty is available)
\IfFileExists{upquote.sty}{\usepackage{upquote}}{}


% ============================================================
% PAGE GEOMETRY
% ============================================================

% Wide left margin (5 cm) for binding; standard margins elsewhere
\usepackage[lmargin=5cm,rmargin=2.5cm,tmargin=2.5cm,bmargin=2.5cm]{geometry}


% ============================================================
% FIGURE PLACEMENT
% ============================================================

% Force all figures to appear at their exact position in the source [H]
\usepackage{float}
\let\origfigure\figure
\let\endorigfigure\endfigure
\renewenvironment{figure}[1][2]{
    \expandafter\origfigure\expandafter[H]
}{
    \endorigfigure
}


% ============================================================
% BIBLIOGRAPHY
% ============================================================

% APA citation style via biblatex + biber backend
% maxbibnames=99 ensures all authors are listed in the reference list
\usepackage{csquotes}
\usepackage[backend=biber, maxbibnames=99, style=apa]{biblatex}
\setlength\bibitemsep{1.5\itemsep}         % Extra spacing between reference entries
\bibliography{Resources/references.bib}    % Path to the .bib file


% ============================================================
% CODE SYNTAX HIGHLIGHTING
% ============================================================
% Required for displaying formatted code blocks.
% Safe to remove this entire section if your thesis contains no code.

\usepackage{color}
\usepackage{fancyvrb}
\newcommand{\VerbBar}{|}
\newcommand{\VERB}{\Verb[commandchars=\\\{\}]}
\DefineVerbatimEnvironment{Highlighting}{Verbatim}{commandchars=\\\{\}}
\usepackage{framed}
\definecolor{shadecolor}{RGB}{248,248,248}
\newenvironment{Shaded}{\begin{snugshade}}{\end{snugshade}}
% Token color definitions for syntax highlighting
\newcommand{\AlertTok}[1]{\textcolor[rgb]{0.94,0.16,0.16}{#1}}
\newcommand{\AnnotationTok}[1]{\textcolor[rgb]{0.56,0.35,0.01}{\textbf{\textit{#1}}}}
\newcommand{\AttributeTok}[1]{\textcolor[rgb]{0.13,0.29,0.53}{#1}}
\newcommand{\BaseNTok}[1]{\textcolor[rgb]{0.00,0.00,0.81}{#1}}
\newcommand{\BuiltInTok}[1]{#1}
\newcommand{\CharTok}[1]{\textcolor[rgb]{0.31,0.60,0.02}{#1}}
\newcommand{\CommentTok}[1]{\textcolor[rgb]{0.56,0.35,0.01}{\textit{#1}}}
\newcommand{\CommentVarTok}[1]{\textcolor[rgb]{0.56,0.35,0.01}{\textbf{\textit{#1}}}}
\newcommand{\ConstantTok}[1]{\textcolor[rgb]{0.56,0.35,0.01}{#1}}
\newcommand{\ControlFlowTok}[1]{\textcolor[rgb]{0.13,0.29,0.53}{\textbf{#1}}}
\newcommand{\DataTypeTok}[1]{\textcolor[rgb]{0.13,0.29,0.53}{#1}}
\newcommand{\DecValTok}[1]{\textcolor[rgb]{0.00,0.00,0.81}{#1}}
\newcommand{\DocumentationTok}[1]{\textcolor[rgb]{0.56,0.35,0.01}{\textbf{\textit{#1}}}}
\newcommand{\ErrorTok}[1]{\textcolor[rgb]{0.64,0.00,0.00}{\textbf{#1}}}
\newcommand{\ExtensionTok}[1]{#1}
\newcommand{\FloatTok}[1]{\textcolor[rgb]{0.00,0.00,0.81}{#1}}
\newcommand{\FunctionTok}[1]{\textcolor[rgb]{0.13,0.29,0.53}{\textbf{#1}}}
\newcommand{\ImportTok}[1]{#1}
\newcommand{\InformationTok}[1]{\textcolor[rgb]{0.56,0.35,0.01}{\textbf{\textit{#1}}}}
\newcommand{\KeywordTok}[1]{\textcolor[rgb]{0.13,0.29,0.53}{\textbf{#1}}}
\newcommand{\NormalTok}[1]{#1}
\newcommand{\OperatorTok}[1]{\textcolor[rgb]{0.81,0.36,0.00}{\textbf{#1}}}
\newcommand{\OtherTok}[1]{\textcolor[rgb]{0.56,0.35,0.01}{#1}}
\newcommand{\PreprocessorTok}[1]{\textcolor[rgb]{0.56,0.35,0.01}{\textit{#1}}}
\newcommand{\RegionMarkerTok}[1]{#1}
\newcommand{\SpecialCharTok}[1]{\textcolor[rgb]{0.81,0.36,0.00}{\textbf{#1}}}
\newcommand{\SpecialStringTok}[1]{\textcolor[rgb]{0.31,0.60,0.02}{#1}}
\newcommand{\StringTok}[1]{\textcolor[rgb]{0.31,0.60,0.02}{#1}}
\newcommand{\VariableTok}[1]{\textcolor[rgb]{0.00,0.00,0.00}{#1}}
\newcommand{\VerbatimStringTok}[1]{\textcolor[rgb]{0.31,0.60,0.02}{#1}}
\newcommand{\WarningTok}[1]{\textcolor[rgb]{0.56,0.35,0.01}{\textbf{\textit{#1}}}}


% ============================================================
% TABLES & GRAPHICS
% ============================================================

\usepackage{longtable,booktabs}  % Multi-page tables and publication-quality rules
\usepackage{graphicx}

% Auto-scale images so they never exceed the text area
\makeatletter
\def\maxwidth{\ifdim\Gin@nat@width>\linewidth\linewidth\else\Gin@nat@width\fi}
\def\maxheight{\ifdim\Gin@nat@height>\textheight\textheight\else\Gin@nat@height\fi}
\makeatother
\setkeys{Gin}{width=\maxwidth,height=\maxheight,keepaspectratio}

% Passthrough wrapper used by figure includes; keeps natural image size
\newcommand{\pandocbounded}[1]{#1}


% ============================================================
% HYPERLINKS & PDF METADATA
% ============================================================

% Update pdfauthor and pdftitle to match your thesis
\ifxetex
  \usepackage[setpagesize=false, % page size defined by xetex
              unicode=false,     % unicode breaks when used with xetex
              xetex]{hyperref}
\else
  \usepackage[unicode=true]{hyperref}
\fi
\hypersetup{
  breaklinks  = true,
  bookmarks   = true,
  pdfauthor   = {Name},   % <-- Your name
  pdftitle    = {Title},  % <-- Your thesis title
  colorlinks  = true,
  citecolor   = blue,
  urlcolor    = blue,
  linkcolor   = black,
  pdfborder   = {0 0 0}
}
\urlstyle{same}  % URLs use the same font as surrounding text


% ============================================================
% PARAGRAPH & SECTION FORMATTING
% ============================================================

\setlength{\parindent}{0pt}           % No paragraph indentation
\setlength{\parskip}{6pt plus 2pt minus 1pt}  % Vertical space between paragraphs
\setlength{\emergencystretch}{3em}    % Prevents overfull lines
\setcounter{secnumdepth}{5}           % Number sections down to subsubparagraph level


% ============================================================
% TITLE PAGE COMMANDS
% ============================================================

\usepackage{titling}

% \subtitle{} command — used on the title page
\newcommand{\subtitle}[1]{
  \posttitle{
    \begin{center}\large#1\end{center}
  }
}

% Title metadata — update these for your thesis
\setlength{\droptitle}{-2em}
  \title{Title}                          % <-- Thesis title
  \pretitle{\vspace{\droptitle}\centering\huge}
  \posttitle{\par}
\subtitle{Subtitle}                      % <-- Thesis subtitle
  \author{Name}                          % <-- Your name
  \preauthor{\centering\large\emph}
  \postauthor{\par}
  \predate{\centering\large\emph}
  \postdate{\par}
  \date{today}                           % <-- Submission date


% ============================================================
% ACRONYMS / LIST OF ABBREVIATIONS
% ============================================================

% glossaries-package in acronym-only mode (no page numbers in list)
\usepackage[acronym,nomain,nonumberlist]{glossaries}
\usepackage{supertabular}
\setacronymstyle{short-long}  % First use: SHORT (Long Form); later: SHORT
\makenoidxglossaries

% Styling: bold acronym names, no blank lines between groups
\renewcommand{\glsnamefont}[1]{\textbf{#1}}
\renewcommand{\glsgroupskip}{}
\setlength{\glsdescwidth}{0.8\linewidth}

% Two-column table style: "Abbreviation | Definition" with header and rule
\newglossarystyle{mysuper}{%
  \setglossarystyle{super}
  \renewcommand{\glossentry}[2]{%
    \glsentryitem{##1}\glstarget{##1}{\glossentryname{##1}} &
    \glossentrydesc{##1}\glspostdescription\space ##2\\[0.75ex]
  }%
  \renewcommand{\superheader}{\textbf{Abbreviation} & \textbf{Definition}\\\hline\endhead}%
}

% Acronym definitions — add your own with: \newacronym{label}{SHORT}{Long Form}
\newacronym{itt}{ITT}{Intention-to-Treat}
\newacronym{iv}{IV}{Instrumental Variables}
\newacronym{sbcf}{SBCF}{Shrinkage Bayesian Causal Forest}
\newacronym{sbcf-iv}{SBCF-IV}{Shrinkage Bayesian Instrumental Variable Causal Forest}


% ============================================================
% LINE SPACING
% ============================================================

\usepackage{setspace}
\onehalfspacing   % 1.5 line spacing (standard requirement for academic work)


% ============================================================
% LANGUAGE — babel loaded based on \thesislanguage (set at top)
% ============================================================
% \IfLanguageName{english}{...}{...} in ude_thesis.tex body
% automatically switches text once babel is loaded here.

\usepackage{ifthen}
\usepackage{iflang}
\usepackage[super]{nth}  % Ordinal suffixes: 1st, 2nd, 3rd ...

\ifthenelse{\equal{\thesislanguage}{german}}{
  \usepackage[ngerman]{babel}
}{
  \usepackage[english]{babel}
}


% ============================================================
% MICROTYPOGRAPHY
% ============================================================

% Improves text justification and glyph spacing where available
\IfFileExists{microtype.sty}{%
  \usepackage[nopatch=footnote]{microtype}
  \UseMicrotypeSet[protrusion]{basicmath}  % Disable protrusion for tt fonts
}{}

% Do NOT compile this file directly — compile ude_thesis.tex.
% ============================================================

\documentclass[11pt,a4paper]{article}


% ============================================================
% LANGUAGE — change this one value to switch the whole document
% Options: english | german
% ============================================================
\newcommand{\thesislanguage}{english}  % <-- Set language here: english | german

% Include statutory declaration page at end of document?
% true  = declaration page printed (required for final submission)
% false = declaration page omitted  (useful for drafts / seminar papers)
\newcommand{\includestatutorydeclaration}{true}  % <-- true | false


% ============================================================
% FONTS & ENCODING
% ============================================================

\usepackage{lmodern}          % Latin Modern font (scalable)
\usepackage{amssymb,amsmath}  % AMS math symbols and environments
\usepackage{ifxetex,ifluatex} % Detect compiler (pdflatex vs xetex/luatex)
\usepackage{fixltx2e}         % Provides \textsubscript (legacy fix)

% Font encoding: T1 for pdflatex, fontspec for xetex/luatex
\ifnum 0\ifxetex 1\fi\ifluatex 1\fi=0
  % pdflatex path
  \usepackage[T1]{fontenc}
  \usepackage[utf8]{inputenc}
\else
  % xetex / luatex path
  \ifxetex
    \usepackage{mathspec}
    \usepackage{xltxtra,xunicode}
  \else
    \usepackage{fontspec}
  \fi
  \defaultfontfeatures{Mapping=tex-text,Scale=MatchLowercase}
  \newcommand{\euro}{€}
\fi

% Straight quotes in verbatim environments (if upquote.sty is available)
\IfFileExists{upquote.sty}{\usepackage{upquote}}{}


% ============================================================
% PAGE GEOMETRY
% ============================================================

% Wide left margin (5 cm) for binding; standard margins elsewhere
\usepackage[lmargin=5cm,rmargin=2.5cm,tmargin=2.5cm,bmargin=2.5cm]{geometry}


% ============================================================
% FIGURE PLACEMENT
% ============================================================

% Force all figures to appear at their exact position in the source [H]
\usepackage{float}
\let\origfigure\figure
\let\endorigfigure\endfigure
\renewenvironment{figure}[1][2]{
    \expandafter\origfigure\expandafter[H]
}{
    \endorigfigure
}


% ============================================================
% BIBLIOGRAPHY
% ============================================================

% APA citation style via biblatex + biber backend
% maxbibnames=99 ensures all authors are listed in the reference list
\usepackage{csquotes}
\usepackage[backend=biber, maxbibnames=99, style=apa]{biblatex}
\setlength\bibitemsep{1.5\itemsep}         % Extra spacing between reference entries
\bibliography{Resources/references.bib}    % Path to the .bib file


% ============================================================
% CODE SYNTAX HIGHLIGHTING
% ============================================================
% Required for displaying formatted code blocks.
% Safe to remove this entire section if your thesis contains no code.

\usepackage{color}
\usepackage{fancyvrb}
\newcommand{\VerbBar}{|}
\newcommand{\VERB}{\Verb[commandchars=\\\{\}]}
\DefineVerbatimEnvironment{Highlighting}{Verbatim}{commandchars=\\\{\}}
\usepackage{framed}
\definecolor{shadecolor}{RGB}{248,248,248}
\newenvironment{Shaded}{\begin{snugshade}}{\end{snugshade}}
% Token color definitions for syntax highlighting
\newcommand{\AlertTok}[1]{\textcolor[rgb]{0.94,0.16,0.16}{#1}}
\newcommand{\AnnotationTok}[1]{\textcolor[rgb]{0.56,0.35,0.01}{\textbf{\textit{#1}}}}
\newcommand{\AttributeTok}[1]{\textcolor[rgb]{0.13,0.29,0.53}{#1}}
\newcommand{\BaseNTok}[1]{\textcolor[rgb]{0.00,0.00,0.81}{#1}}
\newcommand{\BuiltInTok}[1]{#1}
\newcommand{\CharTok}[1]{\textcolor[rgb]{0.31,0.60,0.02}{#1}}
\newcommand{\CommentTok}[1]{\textcolor[rgb]{0.56,0.35,0.01}{\textit{#1}}}
\newcommand{\CommentVarTok}[1]{\textcolor[rgb]{0.56,0.35,0.01}{\textbf{\textit{#1}}}}
\newcommand{\ConstantTok}[1]{\textcolor[rgb]{0.56,0.35,0.01}{#1}}
\newcommand{\ControlFlowTok}[1]{\textcolor[rgb]{0.13,0.29,0.53}{\textbf{#1}}}
\newcommand{\DataTypeTok}[1]{\textcolor[rgb]{0.13,0.29,0.53}{#1}}
\newcommand{\DecValTok}[1]{\textcolor[rgb]{0.00,0.00,0.81}{#1}}
\newcommand{\DocumentationTok}[1]{\textcolor[rgb]{0.56,0.35,0.01}{\textbf{\textit{#1}}}}
\newcommand{\ErrorTok}[1]{\textcolor[rgb]{0.64,0.00,0.00}{\textbf{#1}}}
\newcommand{\ExtensionTok}[1]{#1}
\newcommand{\FloatTok}[1]{\textcolor[rgb]{0.00,0.00,0.81}{#1}}
\newcommand{\FunctionTok}[1]{\textcolor[rgb]{0.13,0.29,0.53}{\textbf{#1}}}
\newcommand{\ImportTok}[1]{#1}
\newcommand{\InformationTok}[1]{\textcolor[rgb]{0.56,0.35,0.01}{\textbf{\textit{#1}}}}
\newcommand{\KeywordTok}[1]{\textcolor[rgb]{0.13,0.29,0.53}{\textbf{#1}}}
\newcommand{\NormalTok}[1]{#1}
\newcommand{\OperatorTok}[1]{\textcolor[rgb]{0.81,0.36,0.00}{\textbf{#1}}}
\newcommand{\OtherTok}[1]{\textcolor[rgb]{0.56,0.35,0.01}{#1}}
\newcommand{\PreprocessorTok}[1]{\textcolor[rgb]{0.56,0.35,0.01}{\textit{#1}}}
\newcommand{\RegionMarkerTok}[1]{#1}
\newcommand{\SpecialCharTok}[1]{\textcolor[rgb]{0.81,0.36,0.00}{\textbf{#1}}}
\newcommand{\SpecialStringTok}[1]{\textcolor[rgb]{0.31,0.60,0.02}{#1}}
\newcommand{\StringTok}[1]{\textcolor[rgb]{0.31,0.60,0.02}{#1}}
\newcommand{\VariableTok}[1]{\textcolor[rgb]{0.00,0.00,0.00}{#1}}
\newcommand{\VerbatimStringTok}[1]{\textcolor[rgb]{0.31,0.60,0.02}{#1}}
\newcommand{\WarningTok}[1]{\textcolor[rgb]{0.56,0.35,0.01}{\textbf{\textit{#1}}}}


% ============================================================
% TABLES & GRAPHICS
% ============================================================

\usepackage{longtable,booktabs}  % Multi-page tables and publication-quality rules
\usepackage{graphicx}

% Auto-scale images so they never exceed the text area
\makeatletter
\def\maxwidth{\ifdim\Gin@nat@width>\linewidth\linewidth\else\Gin@nat@width\fi}
\def\maxheight{\ifdim\Gin@nat@height>\textheight\textheight\else\Gin@nat@height\fi}
\makeatother
\setkeys{Gin}{width=\maxwidth,height=\maxheight,keepaspectratio}

% Passthrough wrapper used by figure includes; keeps natural image size
\newcommand{\pandocbounded}[1]{#1}


% ============================================================
% HYPERLINKS & PDF METADATA
% ============================================================

% Update pdfauthor and pdftitle to match your thesis
\ifxetex
  \usepackage[setpagesize=false, % page size defined by xetex
              unicode=false,     % unicode breaks when used with xetex
              xetex]{hyperref}
\else
  \usepackage[unicode=true]{hyperref}
\fi
\hypersetup{
  breaklinks  = true,
  bookmarks   = true,
  pdfauthor   = {Name},   % <-- Your name
  pdftitle    = {Title},  % <-- Your thesis title
  colorlinks  = true,
  citecolor   = blue,
  urlcolor    = blue,
  linkcolor   = black,
  pdfborder   = {0 0 0}
}
\urlstyle{same}  % URLs use the same font as surrounding text


% ============================================================
% PARAGRAPH & SECTION FORMATTING
% ============================================================

\setlength{\parindent}{0pt}           % No paragraph indentation
\setlength{\parskip}{6pt plus 2pt minus 1pt}  % Vertical space between paragraphs
\setlength{\emergencystretch}{3em}    % Prevents overfull lines
\setcounter{secnumdepth}{5}           % Number sections down to subsubparagraph level


% ============================================================
% TITLE PAGE COMMANDS
% ============================================================

\usepackage{titling}

% \subtitle{} command — used on the title page
\newcommand{\subtitle}[1]{
  \posttitle{
    \begin{center}\large#1\end{center}
  }
}

% Title metadata — update these for your thesis
\setlength{\droptitle}{-2em}
  \title{Title}                          % <-- Thesis title
  \pretitle{\vspace{\droptitle}\centering\huge}
  \posttitle{\par}
\subtitle{Subtitle}                      % <-- Thesis subtitle
  \author{Name}                          % <-- Your name
  \preauthor{\centering\large\emph}
  \postauthor{\par}
  \predate{\centering\large\emph}
  \postdate{\par}
  \date{today}                           % <-- Submission date


% ============================================================
% ACRONYMS / LIST OF ABBREVIATIONS
% ============================================================

% glossaries-package in acronym-only mode (no page numbers in list)
\usepackage[acronym,nomain,nonumberlist]{glossaries}
\usepackage{supertabular}
\setacronymstyle{short-long}  % First use: SHORT (Long Form); later: SHORT
\makenoidxglossaries

% Styling: bold acronym names, no blank lines between groups
\renewcommand{\glsnamefont}[1]{\textbf{#1}}
\renewcommand{\glsgroupskip}{}
\setlength{\glsdescwidth}{0.8\linewidth}

% Two-column table style: "Abbreviation | Definition" with header and rule
\newglossarystyle{mysuper}{%
  \setglossarystyle{super}
  \renewcommand{\glossentry}[2]{%
    \glsentryitem{##1}\glstarget{##1}{\glossentryname{##1}} &
    \glossentrydesc{##1}\glspostdescription\space ##2\\[0.75ex]
  }%
  \renewcommand{\superheader}{\textbf{Abbreviation} & \textbf{Definition}\\\hline\endhead}%
}

% Acronym definitions — add your own with: \newacronym{label}{SHORT}{Long Form}
\newacronym{itt}{ITT}{Intention-to-Treat}
\newacronym{iv}{IV}{Instrumental Variables}
\newacronym{sbcf}{SBCF}{Shrinkage Bayesian Causal Forest}
\newacronym{sbcf-iv}{SBCF-IV}{Shrinkage Bayesian Instrumental Variable Causal Forest}


% ============================================================
% LINE SPACING
% ============================================================

\usepackage{setspace}
\onehalfspacing   % 1.5 line spacing (standard requirement for academic work)


% ============================================================
% LANGUAGE — babel loaded based on \thesislanguage (set at top)
% ============================================================
% \IfLanguageName{english}{...}{...} in ude_thesis.tex body
% automatically switches text once babel is loaded here.

\usepackage{ifthen}
\usepackage{iflang}
\usepackage[super]{nth}  % Ordinal suffixes: 1st, 2nd, 3rd ...

\ifthenelse{\equal{\thesislanguage}{german}}{
  \usepackage[ngerman]{babel}
}{
  \usepackage[english]{babel}
}


% ============================================================
% MICROTYPOGRAPHY
% ============================================================

% Improves text justification and glyph spacing where available
\IfFileExists{microtype.sty}{%
  \usepackage[nopatch=footnote]{microtype}
  \UseMicrotypeSet[protrusion]{basicmath}  % Disable protrusion for tt fonts
}{}

% Do NOT compile this file directly — compile ude_thesis.tex.
% ============================================================

\documentclass[11pt,a4paper]{article}


% ============================================================
% LANGUAGE — change this one value to switch the whole document
% Options: english | german
% ============================================================
\newcommand{\thesislanguage}{english}  % <-- Set language here: english | german

% Include statutory declaration page at end of document?
% true  = declaration page printed (required for final submission)
% false = declaration page omitted  (useful for drafts / seminar papers)
\newcommand{\includestatutorydeclaration}{true}  % <-- true | false


% ============================================================
% FONTS & ENCODING
% ============================================================

\usepackage{lmodern}          % Latin Modern font (scalable)
\usepackage{amssymb,amsmath}  % AMS math symbols and environments
\usepackage{ifxetex,ifluatex} % Detect compiler (pdflatex vs xetex/luatex)
\usepackage{fixltx2e}         % Provides \textsubscript (legacy fix)

% Font encoding: T1 for pdflatex, fontspec for xetex/luatex
\ifnum 0\ifxetex 1\fi\ifluatex 1\fi=0
  % pdflatex path
  \usepackage[T1]{fontenc}
  \usepackage[utf8]{inputenc}
\else
  % xetex / luatex path
  \ifxetex
    \usepackage{mathspec}
    \usepackage{xltxtra,xunicode}
  \else
    \usepackage{fontspec}
  \fi
  \defaultfontfeatures{Mapping=tex-text,Scale=MatchLowercase}
  \newcommand{\euro}{€}
\fi

% Straight quotes in verbatim environments (if upquote.sty is available)
\IfFileExists{upquote.sty}{\usepackage{upquote}}{}


% ============================================================
% PAGE GEOMETRY
% ============================================================

% Wide left margin (5 cm) for binding; standard margins elsewhere
\usepackage[lmargin=5cm,rmargin=2.5cm,tmargin=2.5cm,bmargin=2.5cm]{geometry}


% ============================================================
% FIGURE PLACEMENT
% ============================================================

% Force all figures to appear at their exact position in the source [H]
\usepackage{float}
\let\origfigure\figure
\let\endorigfigure\endfigure
\renewenvironment{figure}[1][2]{
    \expandafter\origfigure\expandafter[H]
}{
    \endorigfigure
}


% ============================================================
% BIBLIOGRAPHY
% ============================================================

% APA citation style via biblatex + biber backend
% maxbibnames=99 ensures all authors are listed in the reference list
\usepackage{csquotes}
\usepackage[backend=biber, maxbibnames=99, style=apa]{biblatex}
\setlength\bibitemsep{1.5\itemsep}         % Extra spacing between reference entries
\bibliography{Resources/references.bib}    % Path to the .bib file


% ============================================================
% CODE SYNTAX HIGHLIGHTING
% ============================================================
% Required for displaying formatted code blocks.
% Safe to remove this entire section if your thesis contains no code.

\usepackage{color}
\usepackage{fancyvrb}
\newcommand{\VerbBar}{|}
\newcommand{\VERB}{\Verb[commandchars=\\\{\}]}
\DefineVerbatimEnvironment{Highlighting}{Verbatim}{commandchars=\\\{\}}
\usepackage{framed}
\definecolor{shadecolor}{RGB}{248,248,248}
\newenvironment{Shaded}{\begin{snugshade}}{\end{snugshade}}
% Token color definitions for syntax highlighting
\newcommand{\AlertTok}[1]{\textcolor[rgb]{0.94,0.16,0.16}{#1}}
\newcommand{\AnnotationTok}[1]{\textcolor[rgb]{0.56,0.35,0.01}{\textbf{\textit{#1}}}}
\newcommand{\AttributeTok}[1]{\textcolor[rgb]{0.13,0.29,0.53}{#1}}
\newcommand{\BaseNTok}[1]{\textcolor[rgb]{0.00,0.00,0.81}{#1}}
\newcommand{\BuiltInTok}[1]{#1}
\newcommand{\CharTok}[1]{\textcolor[rgb]{0.31,0.60,0.02}{#1}}
\newcommand{\CommentTok}[1]{\textcolor[rgb]{0.56,0.35,0.01}{\textit{#1}}}
\newcommand{\CommentVarTok}[1]{\textcolor[rgb]{0.56,0.35,0.01}{\textbf{\textit{#1}}}}
\newcommand{\ConstantTok}[1]{\textcolor[rgb]{0.56,0.35,0.01}{#1}}
\newcommand{\ControlFlowTok}[1]{\textcolor[rgb]{0.13,0.29,0.53}{\textbf{#1}}}
\newcommand{\DataTypeTok}[1]{\textcolor[rgb]{0.13,0.29,0.53}{#1}}
\newcommand{\DecValTok}[1]{\textcolor[rgb]{0.00,0.00,0.81}{#1}}
\newcommand{\DocumentationTok}[1]{\textcolor[rgb]{0.56,0.35,0.01}{\textbf{\textit{#1}}}}
\newcommand{\ErrorTok}[1]{\textcolor[rgb]{0.64,0.00,0.00}{\textbf{#1}}}
\newcommand{\ExtensionTok}[1]{#1}
\newcommand{\FloatTok}[1]{\textcolor[rgb]{0.00,0.00,0.81}{#1}}
\newcommand{\FunctionTok}[1]{\textcolor[rgb]{0.13,0.29,0.53}{\textbf{#1}}}
\newcommand{\ImportTok}[1]{#1}
\newcommand{\InformationTok}[1]{\textcolor[rgb]{0.56,0.35,0.01}{\textbf{\textit{#1}}}}
\newcommand{\KeywordTok}[1]{\textcolor[rgb]{0.13,0.29,0.53}{\textbf{#1}}}
\newcommand{\NormalTok}[1]{#1}
\newcommand{\OperatorTok}[1]{\textcolor[rgb]{0.81,0.36,0.00}{\textbf{#1}}}
\newcommand{\OtherTok}[1]{\textcolor[rgb]{0.56,0.35,0.01}{#1}}
\newcommand{\PreprocessorTok}[1]{\textcolor[rgb]{0.56,0.35,0.01}{\textit{#1}}}
\newcommand{\RegionMarkerTok}[1]{#1}
\newcommand{\SpecialCharTok}[1]{\textcolor[rgb]{0.81,0.36,0.00}{\textbf{#1}}}
\newcommand{\SpecialStringTok}[1]{\textcolor[rgb]{0.31,0.60,0.02}{#1}}
\newcommand{\StringTok}[1]{\textcolor[rgb]{0.31,0.60,0.02}{#1}}
\newcommand{\VariableTok}[1]{\textcolor[rgb]{0.00,0.00,0.00}{#1}}
\newcommand{\VerbatimStringTok}[1]{\textcolor[rgb]{0.31,0.60,0.02}{#1}}
\newcommand{\WarningTok}[1]{\textcolor[rgb]{0.56,0.35,0.01}{\textbf{\textit{#1}}}}


% ============================================================
% TABLES & GRAPHICS
% ============================================================

\usepackage{longtable,booktabs}  % Multi-page tables and publication-quality rules
\usepackage{graphicx}

% Auto-scale images so they never exceed the text area
\makeatletter
\def\maxwidth{\ifdim\Gin@nat@width>\linewidth\linewidth\else\Gin@nat@width\fi}
\def\maxheight{\ifdim\Gin@nat@height>\textheight\textheight\else\Gin@nat@height\fi}
\makeatother
\setkeys{Gin}{width=\maxwidth,height=\maxheight,keepaspectratio}

% Passthrough wrapper used by figure includes; keeps natural image size
\newcommand{\pandocbounded}[1]{#1}


% ============================================================
% HYPERLINKS & PDF METADATA
% ============================================================

% Update pdfauthor and pdftitle to match your thesis
\ifxetex
  \usepackage[setpagesize=false, % page size defined by xetex
              unicode=false,     % unicode breaks when used with xetex
              xetex]{hyperref}
\else
  \usepackage[unicode=true]{hyperref}
\fi
\hypersetup{
  breaklinks  = true,
  bookmarks   = true,
  pdfauthor   = {Name},   % <-- Your name
  pdftitle    = {Title},  % <-- Your thesis title
  colorlinks  = true,
  citecolor   = blue,
  urlcolor    = blue,
  linkcolor   = black,
  pdfborder   = {0 0 0}
}
\urlstyle{same}  % URLs use the same font as surrounding text


% ============================================================
% PARAGRAPH & SECTION FORMATTING
% ============================================================

\setlength{\parindent}{0pt}           % No paragraph indentation
\setlength{\parskip}{6pt plus 2pt minus 1pt}  % Vertical space between paragraphs
\setlength{\emergencystretch}{3em}    % Prevents overfull lines
\setcounter{secnumdepth}{5}           % Number sections down to subsubparagraph level


% ============================================================
% TITLE PAGE COMMANDS
% ============================================================

\usepackage{titling}

% \subtitle{} command — used on the title page
\newcommand{\subtitle}[1]{
  \posttitle{
    \begin{center}\large#1\end{center}
  }
}

% Title metadata — update these for your thesis
\setlength{\droptitle}{-2em}
  \title{Title}                          % <-- Thesis title
  \pretitle{\vspace{\droptitle}\centering\huge}
  \posttitle{\par}
\subtitle{Subtitle}                      % <-- Thesis subtitle
  \author{Name}                          % <-- Your name
  \preauthor{\centering\large\emph}
  \postauthor{\par}
  \predate{\centering\large\emph}
  \postdate{\par}
  \date{today}                           % <-- Submission date


% ============================================================
% ACRONYMS / LIST OF ABBREVIATIONS
% ============================================================

% glossaries-package in acronym-only mode (no page numbers in list)
\usepackage[acronym,nomain,nonumberlist]{glossaries}
\usepackage{supertabular}
\setacronymstyle{short-long}  % First use: SHORT (Long Form); later: SHORT
\makenoidxglossaries

% Styling: bold acronym names, no blank lines between groups
\renewcommand{\glsnamefont}[1]{\textbf{#1}}
\renewcommand{\glsgroupskip}{}
\setlength{\glsdescwidth}{0.8\linewidth}

% Two-column table style: "Abbreviation | Definition" with header and rule
\newglossarystyle{mysuper}{%
  \setglossarystyle{super}
  \renewcommand{\glossentry}[2]{%
    \glsentryitem{##1}\glstarget{##1}{\glossentryname{##1}} &
    \glossentrydesc{##1}\glspostdescription\space ##2\\[0.75ex]
  }%
  \renewcommand{\superheader}{\textbf{Abbreviation} & \textbf{Definition}\\\hline\endhead}%
}

% Acronym definitions — add your own with: \newacronym{label}{SHORT}{Long Form}
\newacronym{itt}{ITT}{Intention-to-Treat}
\newacronym{iv}{IV}{Instrumental Variables}
\newacronym{sbcf}{SBCF}{Shrinkage Bayesian Causal Forest}
\newacronym{sbcf-iv}{SBCF-IV}{Shrinkage Bayesian Instrumental Variable Causal Forest}


% ============================================================
% LINE SPACING
% ============================================================

\usepackage{setspace}
\onehalfspacing   % 1.5 line spacing (standard requirement for academic work)


% ============================================================
% LANGUAGE — babel loaded based on \thesislanguage (set at top)
% ============================================================
% \IfLanguageName{english}{...}{...} in ude_thesis.tex body
% automatically switches text once babel is loaded here.

\usepackage{ifthen}
\usepackage{iflang}
\usepackage[super]{nth}  % Ordinal suffixes: 1st, 2nd, 3rd ...

\ifthenelse{\equal{\thesislanguage}{german}}{
  \usepackage[ngerman]{babel}
}{
  \usepackage[english]{babel}
}


% ============================================================
% MICROTYPOGRAPHY
% ============================================================

% Improves text justification and glyph spacing where available
\IfFileExists{microtype.sty}{%
  \usepackage[nopatch=footnote]{microtype}
  \UseMicrotypeSet[protrusion]{basicmath}  % Disable protrusion for tt fonts
}{}
